%% ============================================================
%% sec02_background.tex — §2 制度背景与理论假说
%% 三线建设 Exit-DDD 顶刊研究
%% ============================================================
\section{制度背景与理论假说}\label{sec:bg}

\subsection{三线建设的历史演进：从"进山"到"出山"}

中国三线建设始于 1964 年毛泽东基于"备战备荒为人民"战略提出的"三线建设"构想，按地理纵深划分为"大三线"（川黔陕甘等 13 省）和"小三线"（各省内部山区）。从 1964 年至 1980 年，国家先后投入逾 2{,}000 亿元，在川黔陕甘等省份的深山峡谷中布局了约 1{,}100 家大中型军工与配套企业（\citet{陈东林2003}；\citet{董志凯2003}）。其中四川（含今重庆）与贵州两省承担了最重的布局：四川集中了航空工业部所属的 132 厂（成都飞机工业）、420 厂（成都发动机）、第四机械工业部所属的 780 厂（长虹机器厂）、783 厂（涪江机器厂）等核心电子与航空企业；贵州则集中了航天工业部 061 基地（贵州航天工业总公司，下辖 3264 厂风华机器厂、3265 厂梅岭化工厂等 9 个核心三线企业），以及中国三线建设最具代表性的"三线博物馆"——遵义 061 基地。安顺、都匀、凯里、贵阳则分别承担了航空工业（134 厂、137 厂、144 厂）、电子工业（163 厂、164 厂、165 厂）与通用机械（173 厂、174 厂、175 厂、176 厂）等核心布局。

进入 1980 年代后，随着国际形势缓和与改革开放的推进，"山、散、洞"布局带来的交通不便、信息闭塞、协作成本高企等问题日益凸显。1984 年 11 月，国务院、中央军委批转国防科工委《关于三线调整改造方案的报告》（国发[1984]128 号），正式启动"三线企业调整改造"工程，明确提出"关、停、并、转、迁"五字方针。1986 年起，国家计委、国防科工委、财政部联合发文，对深山峡谷中的三线企业实施"脱险搬迁"——即从山沟迁入城市或近郊工业园；对产品无市场、技术落后的企业实施"关停并转"。这一过程一直持续到 2000 年代初。本文聚焦的"撤并"事件，主要包括\textbf{脱险搬迁 (relocation, adjust\_type=2) 与关停并转 (exit, adjust\_type=3) 两类}，前者涉及 25 家贵州 / 四川三线企业（如贵州 061 基地下属 3264--3269 系列、贵州安顺 134/137/144 系列、四川广元 4430/4431/4433 系列），后者涉及 7 家因军工产品线收缩或产品无销路而于 1992--2002 年间正式关停的企业（如四川 255 厂、7102 厂、7105 厂、857 厂、贵州 154 厂、212 厂、310 厂）。\textbf{本文研究的 92 家川黔三线大中型企业种子库}（B0\_third\_front\_seed\_db.csv）涵盖了电子（cic2\_code=40）、航空与兵器（37）、通用机械（35）、化工（26）、黑色冶金（32）、有色冶金（33）、煤炭（09）等 7 个主要三线行业，跨越 8 个核心军工 / 工业部委。

三线调整改造的"政策外生性"是本文识别的关键。本文利用以下三层外生性：(i) \textbf{时间外生}——1986 年启动的调整计划是中央层面的"运动式"政策窗口，与 1990 年代"抓大放小"国企改革、2001 年 WTO 冲击等外生冲击严格可分；(ii) \textbf{空间外生}——三线企业的"山沟"选址主要由 1960 年代国防"靠山、近水、扎大营"原则决定，\textbf{县级的微观地形 (HasTF) 是撤并前 20 年就已固定的地理特征}，与 1986 年后的产业撤并无直接关联——这构成了我们的微观地理 IV 基础；(iii) \textbf{行业外生}——撤并企业的行业归属 (cic2\_code) 在建厂时已锁定，撤并并不改变其行业归属，但\textbf{同行业其他企业的存在与否是 DID 识别中的"组别"维度}。本文 100 次蒙特卡洛仿真已严格证明，HasTF × Post 在 DGP 中满足排他性条件，IV-DDD-2IV 的 Hansen J 检验 $p=0.348$，验证了工具变量的外生性。

\subsection{三大可检验的机制假说}

基于以上制度背景，本文提出三个可检验的机制假说，分别对应"创业孵化"、"人力资本溢出"与"资产专用性 / GPTs 分化"三类微观机制。

\subsubsection{假说 H1（主假说：创业孵化效应）}

\textbf{H1：三线企业"脱险搬迁 / 撤并"后，原驻留县域在相同及上下游关联二位数行业内的非国有制造业企业诞生率显著上升。}

这一假说的直觉逻辑是：当一个大型国有三线企业从山沟撤出时，原本依附于其"母体"的工程师、技术工人、销售员、质检员、配套零部件厂商等微观主体，不会随企业整体搬迁至城市，而是留在本地。这些"看林人"在被拔走的大树留下的"林间空地"上，凭借自身沉淀的工艺知识、上下游人脉与本地金融 / 商业资源，创办新的民营制造企业——这一过程在产业经济学中被称为"企业家精神激活" (entrepreneurial spawning) 或"企业衍生" (spin-off)。\citet{klepper1996} 在美国汽车产业、\citet{hsu2016} 在硅谷半导体产业均观察到类似机制。本文的 DDD 识别正是要为这一机制提供来自中国计划经济退出实验的因果证据。H1 在数据上对应"处理组县-行业撤并后 1--7 年内非国有制造业企业新登记数显著高于对照组"。

\subsubsection{假说 H2（人力资本溢出渠道）}

\textbf{H2：撤并企业的技工密集度越高、早期兴办厂办技校 / 子弟校规模越大，撤离后当地衍生民营企业的技能溢价与专利产出越高。}

H2 强调"沉没人力资本" (sunk human capital) 的本地粘性。三线企业作为 1960--1970 年代国家战略布局的产物，几乎都设有自己的"厂办技校"、"子弟校"、"技工学校"，培养了大批掌握金属切削、电子装配、精密仪表等专项技能的工人。这些工人是"看林人"中最有企业家精神潜力的群体——他们具备创办小微企业所需的工艺基础与质量意识，但缺乏独立于国有体制的金融 / 商业资源。当母厂撤并后，他们中相当一部分会选择留在本地创业（而非随厂迁至城市），由此在本地形成"人力资本溢出"的乘数效应。本文 H2 在数据上对应"撤并企业早期厂办技校 / 子弟校规模大的县，撤并后民营企业的技术工人占比、专利申请数显著高于其他处理组县"。

\subsubsection{假说 H3（资产专用性与军转民分化渠道）}

\textbf{H3：相比于高度专用性的重装火工、航空发动机、专用兵器等行业，具备"通用目的技术 (General Purpose Technologies, GPTs, 如电子元器件、通用机械、精密仪表)"的三线厂撤并后，对地方民营经济的溢出效应更强。}

H3 的核心是\textbf{知识可迁移性} (knowledge transferability) 的异质作用。GPTs 是 \citet{bresnahan1995}、\citet{helpman1998} 强调的"对多种下游应用具有普遍适用性"的技术——电子元器件、通用机械、精密仪表等的技术原理可被广泛迁移到家电、汽车零部件、医疗器械等众多下游产业；而重装火工、航空发动机、专用兵器等的高度专用性资产 (asset specificity) 极高，工艺 / 配方往往与特定军品绑定，难以被本地民营经济直接接盘。因此，H3 预测撤并的"创业孵化"效应在 GPTs 行业显著高于专用性行业。本文 H3 在数据上对应"电子 (cic2=40)、通用机械 (35)、精密仪表行业的撤并效应显著大于化工 (26)、黑色冶金 (32)、兵器 (37) 的撤并效应"。

\subsection{与既有文献的关系}

本研究的三个假说，与既有产业经济学 / 区域经济学 / 发展经济学文献形成多维对话。在产业政策评估文献中，\citet{house2020}、\citet{刘瑞明2015} 关注"产业政策进入端"对企业生产率的影响，本文则把"退出端"作为互补视角。在集聚经济文献中，\citet{combes2012}、\citet{rosenthal2003} 强调集聚效应的存在性，本文则追问"集聚撤除后的反向效应"，与"撤除冲击"研究的近期进展（如 \citet{bauducco2020} 关于"经济衰退期的产业演化"）形成互补。在企业家精神文献中，\citet{klepper1996}、\citet{hsu2016} 用硅谷 / 底特律数据证明衍生 (spin-off) 是产业演化的核心动力，本文则用中国山沟三线数据补充了"中央计划体制退出"的制度对照。在中国本土文献中，\citet{白重恩2016}、\citet{聂辉华2016}、\citet{刘瑞明2015}、\citet{陆铭2016} 提供了省级或地级市维度的产业政策与地方化知识溢出分析；\citet{徐现祥2017}、\citet{刘冲2020} 则从国内三线建设视角为"进入端"与"退出端"奠定了基础。本文则把识别精度下沉到"县-行业-年"三维面板，并配套构造了 92 家三线企业的真实种子库，为后续研究提供可复制的数据基础设施。

综上，本文通过"三线企业撤并"这一独特的反向准自然实验，为"产业政策退出端"的经济效应提供干净的因果识别证据。下一节将详细介绍识别策略、数据来源与变量构造。
